\section{Persistent Subsystems: Network Closure}

First we learned what can move (protocols), now we learn what can stay (composites), and why everything else must decay.

As established in the preceding section, the fractional routing protocols (quarks) distribute their flux across the 3-dimensional internal color bulk. 

Because the macroscopic spacetime manifold ($D=4$) only accepts integer-quantized topological boundaries ($\chi=2$), fractional routing protocols cannot propagate independently. If a single fractional protocol attempts to cross spacetime, its fractional boundary flux (e.g., $1/3$) presents an unresolvable rounding error to the discrete causal manifold.

When a fractional routing state is paired with its conjugate $1/3 + 2/3 = 1$ or $1/3 + 1/3 + 1/3 = 1$, it forms a \textbf{Closed Phase Loop} within the internal $N_{res} = 3$ interaction bulk. These combinations yield a net topological footprint equal to zero or an integer, allowing them to be projected as observable matter. This is achieved by forming \textbf{Persistent Subsystems}: fully closed informational loops that independently satisfy the six architectural pillars of Informational Energetics:

\begin{equation}
    M_{subsystem} = \text{CAP} + \text{MAP} + \text{PRO} + \text{GOV} + \text{TOL} + \text{MAR}
\end{equation}

Because the $E_8$ substrate is discrete, there is a strictly limited ``menu'' of valid geometric operators available to construct these subsystems. \textbf{Hadronization} is therefore the algorithmic process of \textbf{Network Closure}---the combinatorial binding of fractional routing protocols until their net topological footprint strictly satisfies these six pillars, producing composite particles (Baryons and Mesons). 

\subsection{The Subsystem Architecture (The IE Menu)}

The first four pillars define the \textbf{Static Topological Mass}, the base geometric routing latency of the isolated structure.
\begin{itemize}
    \item \textbf{CAP (The Hardware Envelope):} The base dimensional footprint required to host the internal state. 
    \begin{itemize}
      \item{0D Point} $\Delta^0 = 1$
      \item{1D Line} $\Delta^1$
      \item{2D Area} $\Delta^2$
      \item{3D Volume} $\Delta^3$
    \end{itemize}

    \item \textbf{MAP (Identity/Boundary): The Topological Anchor}
    \begin{itemize}
        \item \textit{Free Anchor:} Bind to the manifold directly ($D$).
        \item \textit{Conjugate Loop:} Bind to the vacuum impedance ($2\alpha^{-1} - \sigma\chi$).
        \item \textit{Closed Surface:} Bind via a continuous flux tube ($-\pi D$).
    \end{itemize}

    \item \textbf{PRO (Protocol): Internal Synchronization}
    \begin{itemize}
        \item \textit{Free Streaming:} Pay the spin phase lag ($-2\pi$) and gauge load.
        \item \textit{Chiral Oscillation:} Flip across the active chiral channels ($+1/(\nu-\chi)$).
        \item \textit{Orthogonal Exhaustion (Confinement):} Perfectly synchronize all 3 internal dimensions so external routing cost drops to exactly ($0$).
    \end{itemize}

    \item \textbf{GOV (Governor): The Capacity Bounds}
    \begin{itemize}
        \item \textit{Unconstrained:} $0$
        \item \textit{Volumetric Compression:} The 4D manifold pushes inward to contain a 3D bulk ($-1/D$).
        \item \textit{Volumetric Leakage:} The 3D bulk overflows the 4D manifold ($+1/D$).
    \end{itemize}
\end{itemize}

The final two pillars define the \textbf{Dynamic Latency}---the continuous entropic cost of operating as a forward-propagating temporal standing wave on an active substrate.
\begin{itemize}
    \item \textbf{TOL (The Entropic Toll):} The dynamic cost of continuous temporal updates ($Z_{TOL}$). A coherent standing wave is not a static object; it must step forward in time. It pays this exact transition cost for every interaction dimension it actively synchronizes.
    \item \textbf{MAR (The Margin):} The static vacuum noise floor ($Z_{MAR}$). Because the substrate naturally hums at this minimal thermal resolution, a persistent standing wave receives this baseline zero-point energy ``for free'', subtracting from its net dynamic latency.
\end{itemize}

\CatchFileBetweenTags{\ProtonVal}{particles.tex}{ProtonVal}
\CatchFileBetweenTags{\ProtonEq}{particles.tex}{ProtonEq}
\CatchFileBetweenTags{\ProtonExperimentalValue}{particles.tex}{ProtonExperimentalValue}
\CatchFileBetweenTags{\ProtonSigma}{particles.tex}{ProtonSigma}

\subsection{Strategy 1: Orthogonal Exhaustion (The Baryons)}

A Baryon (such as the Proton) achieves network closure by exhausting the internal capacity of the lattice. It routes data simultaneously across all three internal color dimensions, perfectly balancing its internal phase. 

However, a Baryon is a \textbf{Coherent Standing Wave}. It is not a static configuration; it requires continuous microscopic updates to maintain coherence as it propagates. Its total mass is therefore its Static Topological Mass plus its Net Dynamic Latency.

\paragraph{The Static Topological Mass:}
\begin{enumerate}
    \item \textbf{CAP ($\Delta^2$):} Requires an areal resonance to host a 3-axis internal state.
    \item \textbf{MAP ($-\pi D$):} Projects a closed circular flux tube onto the 4D manifold.
    \item \textbf{PRO ($0$):} \textit{Color Confinement}. Perfectly balanced internal symmetry yields zero external synchronization cost.
    \item \textbf{GOV ($-1/D$):} The 4D manifold exerts volumetric compression on the 3D internal geometry.
\end{enumerate}

\paragraph{The Net Dynamic Latency (TOL - MAR):}
The rate of temporal updates is set by the particle's internal synchronization requirements, but the vacuum provides a baseline zero-point energy lift (the Margin) that reduces this latency. 
\begin{itemize}
    \item \textbf{The Proton ($\sigma = 5$):} Because the Proton possesses a net topological boundary flux (an electric charge of $+1$), it is anchored to both the macroscopic spacetime boundary ($\chi=2$) and the internal color bulk ($\sigma-\chi=3$). To maintain its standing wave coherence, it must synchronize its updates across all $\mathbf{5}$ interaction dimensions ($+5Z_{TOL}$). 
    
    Crucially, because its charge is anchored to the topological boundary ($\chi=2$), its interaction with the vacuum's zero-point noise floor ($Z_{MAR}$) is geometrically distributed across this boundary (acting as a spinor on the double cover). It therefore receives exactly half the vacuum lift. Its net dynamic cost is exactly $+5Z_{TOL} - Z_{MAR}/\chi$.
    
    \item \textbf{The Neutron ($\sigma - \chi = 3$):} The Neutron is electrically neutral; its boundary flux is geometrically canceled out. It does not need to synchronize a boundary charge. It \textit{only} synchronizes the 3-dimensional internal color bulk ($+3Z_{TOL}$). 
    
    Because it possesses no net boundary flux, it interacts with the vacuum's zero-point energy purely volumetrically as a whole node, receiving the full, undivided integer vacuum lift ($-Z_{MAR}$). Its net dynamic cost is exactly $+3Z_{TOL} - Z_{MAR}$.
\end{itemize}

Summing these ledgers yields the exact bare masses of the Proton and Neutron:
\begin{equation}
    M_{p (\text{bare})} = \left(\Delta^2 - \pi D - \frac{1}{D}\right) + 5Z_{TOL} - \frac{Z_{MAR}}{\chi}
\end{equation}
\begin{equation}
    M_{n (\text{bare})} = \left(\Delta^2 - \pi D - \frac{1}{D} + \frac{\sigma}{2}\right) + 3Z_{TOL} - Z_{MAR}
\end{equation}

\paragraph{The Strong CP Solution:}
To achieve Orthogonal Exhaustion, the Baryon must perfectly balance its routing across all 3 internal color dimensions. Any CP-violating phase ($\theta_{QCD} \neq 0$) would break this perfect symmetry, reintroducing a positive Protocol latency. This would cause the continuous thermodynamic drag to exceed the capacity limit, triggering immediate decay. Therefore, the infinite stability of the Proton \textit{strictly requires} $\theta_{QCD} = 0$. 

\subsection{Strategy 2: The Conjugate Loop (The Mesons)}

A Meson (such as the Pions and Kaons) achieves stability not by exhausting the internal dimensions, but by pairing a fractional routing protocol with its anti-chiral inverse to form a closed loop. 

Because Mesons are composed of a particle/anti-particle conjugate pair, they possess conjugate temporal phases ($+t$ and $-t$). Internally, their causal flow strictly cancels out. Because they do not form a macroscopic, forward-propagating temporal standing wave, they completely bypass the continuous $Z_{TOL}$ dynamic latency required by the Baryons.

\paragraph{The Pions (Gen 1 Conjugates):}
Pions pair identical generational tiers (e.g., $u\bar{u}$, $d\bar{d}$). 
\begin{enumerate}
    \item \textbf{CAP ($\emptyset$):} The loop requires no macroscopic spatial envelope.
    \item \textbf{MAP ($2\alpha^{-1} - \sigma\chi$):} The identity is a double-load on the baseline vacuum impedance ($2\alpha^{-1}$), adjusted for the structural intersection of the lattice symmetry and topological boundary ($\sigma\chi$).
    \item \textbf{PRO ($+\frac{1}{\nu-\chi}$):} To maintain unanchored resonance, the loop must continuously flip chirality, paying a reciprocal latency cost across the active chiral channels.
\end{enumerate}
Summing the ledger yields the exact bare mass of the neutral Pion: 
\begin{equation}
    M_{\pi^0 (\text{bare})} = 2\alpha^{-1} - \sigma\chi + \frac{1}{\nu-\chi}
\end{equation}

\paragraph{The Kaons (Cross-Generational Conjugates):}
Kaons pair the linear Generation 1 boundary with the areal Generation 2 boundary (the Strange quark). Their base topology expands to $(\sigma+\chi)\alpha^{-1} - \sigma$. The neutral Kaon ($K^0$) adds the down-quark's strong coupling to the 5-fold internal symmetry ($+D\sigma$), while the charged Kaon ($K^+$) adds the up-quark's weak coupling to the circular boundary ($+D\pi$). 

This geometric distinction ($D\sigma$ vs $\pi D$) perfectly explains the famous mass inversion anomaly where the neutral $K^0$ is strictly heavier than the charged $K^+$.

Furthermore, because Kaons cross generations, they must resolve internal transition latencies:
\begin{itemize}
    \item \textbf{The $K^0$ (Isospin-Aligned Loop):} The $K^0$ ($d\bar{s}$) contains two Down-type quarks ($I_3 = -1/2$). Because their isospin is aligned, they do not project across the Weak $W^\pm$ interface. Instead, the loop continuously leaks phase across the generational boundary. This geometric leak is exactly twice the Cabibbo Angle ($2s_{12} = \pi/7$), reducing its latency:
    \begin{equation}
        M_{K^0 (\text{bare})} = (\sigma+\chi)\alpha^{-1} - \sigma + D\sigma - 2\left(\frac{\pi}{\nu-\chi}\right)
    \end{equation}

    \item \textbf{The $K^+$ (Isospin-Crossed Loop):} The $K^+$ ($u\bar{s}$) pairs opposite isospin states ($I_3 = +1/2$ and $-1/2$). To form a coherent loop, they \textit{must} continuously project across the Weak $W^\pm$ interface. This incurs the exact 1D chiral projection factor ($1/\sqrt{2}$) derived in the Protocol architecture.
\end{itemize}


\subsection{Experimental Validation: The Observer Effects (Charge Radiation)}

Evaluating the bare geometric formulas for the neutral Neutron, $\pi^0$, and $K^0$ matches experimental measurements to extreme precision. However, their charged counterparts (the Proton, $\pi^\pm$, and $K^+$) exhibit larger residuals when measured in the continuous macroscopic universe. 

This residual is the \textbf{Macroscopic Charge Load (The EM Tax)}. When a closed subsystem possesses a net boundary flux (Charge), it continuously synchronizes with the surrounding vacuum, radiating a macroscopic electric field. This acts as negative binding energy, reducing the observed mass.

\begin{itemize}
    \item \textbf{Proton (Volumetric Radiation):} As an Orthogonal Exhaustion, the Proton projects a volumetric scalar ($D + 1/D$). Its required charge load is $-\alpha(D + 1/D)$. 
    
    \item \textbf{Charged Mesons (Loop Radiation):} The $\pi^\pm$ and $K^+$ are Conjugate Loops. Their charge radiates from the internal color bulk ($\sigma - \chi$) through the two topological boundaries ($\chi$). Their geometric extent is the strict ratio of these spaces, yielding an identical charge load for both charged mesons: $-\alpha(\frac{\sigma - \chi}{\chi})$. 
\end{itemize}

This strict bifurcation—neutral particles possess zero macroscopic charge load, while charged particles radiate exactly their geometric extent multiplied by $\alpha$—eliminates the need for arbitrary empirical field dressings. Applying this uniform EM Tax to the $K^+$ yields its final observed mass:

\begin{equation}
    M_{K^+ (\text{observed})} = M_{K^+ (\text{bare})} - \frac{1}{\sqrt{2}} - \alpha\left(\frac{\sigma - \chi}{\chi}\right)
\end{equation}



\subsection{Geometric Redundancy (Exotic States)}
This network-closure logic naturally predicts both the existence and the extreme instability of exotic composite states. If Orthogonal Exhaustion (3 parts) and Conjugate Loops (2 parts) achieve network closure, the lattice can theoretically process composite super-structures like \textbf{Tetraquarks} ($qq\bar{q}\bar{q}$) and \textbf{Pentaquarks} ($qqqq\bar{q}$). 

However, because these configurations exceed the \textit{minimum} required bits for topological closure, they introduce massive \textbf{Geometric Redundancy}. By the principles of Informational Energetics, redundant structural complexity generates severe thermodynamic latency. The network violently re-routes these redundant configurations into the lower-latency ground states in yoctoseconds. They exist as valid mathematical solutions to the boundary conditions, but are brutally suppressed by the Survivorship Filter.


\section{The Dynamics of Standing Waves: Extreme Precision Limits}

Evaluating the geometric mass formulas against the CODATA 2022 and PDG 2024 experimental values reveals a striking triumph for the dynamic standing wave architecture. 

\subsection{The Sub-Sigma Baryon Validation}
Previous continuous approximations of the Baryon masses using only static geometric variables ($\Delta^2 - \pi D - 1/D$) yielded an impressive $3$ parts in $10^8$ absolute accuracy, but left an unexplained $\sim 29$ eV residual. 

By correctly modeling the Baryons as \textbf{Coherent Standing Waves} subject to the fundamental transition Toll ($Z_{TOL}$) and the thermal vacuum Margin ($Z_{MAR}$), this residual completely vanishes, confirming the exact physical dynamics of the $E_8$ substrate:

\begin{itemize}
    \item \textbf{The Proton:} By paying the Toll across all $\sigma=5$ interaction dimensions and receiving a boundary-diluted vacuum lift ($Z_{MAR}/\chi$), the dynamic correction adds exactly $+29.54$ eV to the static mass. The calculated mass ($938.27208943$ MeV) matches the CODATA 2022 experimental target ($938.27208943 \pm 0.00000029$ MeV) to within an unprecedented \textbf{0.01$\sigma$}.
    
    \item \textbf{The Neutron:} By paying the Toll strictly across the internal color bulk ($\sigma-\chi=3$) and receiving the full volumetric vacuum lift ($Z_{MAR}$), the dynamic correction adds exactly $+16.68$ eV. The calculated mass ($939.56542245$ MeV) matches the CODATA 2022 experimental target ($939.56542194 \pm 0.00000048$ MeV) to within \textbf{1.06$\sigma$}.
\end{itemize}

This sub-sigma convergence proves that the mass of a Baryon is not a static integer rock, but an active, persistent computational process whose latency is rigidly governed by the exact number of dimensions it must synchronize during each causal update cycle.


\section{ Mass}

In standard model-building, the masses of these fermions are free parameters, dictated by arbitrary Yukawa couplings to the Higgs field. Having established that charge, color, and spin are pure routing mechanics, and the projected subsystems we can compare their energy to experimental results.

TODO import the script equations, experimental values, sigma, etc